\section{Paired-intervention identifiability calculations}
This section is a theoretical extension motivating the next PathoAlign experiments. It is not part of the detector-switch evidence above.

\subsection{Why pairing changes the information structure}
Assume a pathology representation is generated by biological and acquisition factors:
\begin{equation}
  x=Bz_b+Az_a+\varepsilon.
\end{equation}
For a paired acquisition intervention preserving biology,
\begin{align}
  x_i&=Bz_{b,i}+Az_{a,i}+\varepsilon_i,\\
  x_i'&=Bz_{b,i}+Az'_{a,i}+\varepsilon_i'.
\end{align}
Subtracting gives
\begin{equation}
  \Delta x_i=x_i'-x_i=A(z'_{a,i}-z_{a,i})+(\varepsilon_i'-\varepsilon_i),
\end{equation}
so the shared biological term cancels. If acquisition latents have covariance $\Sigma_a$ and the noise is isotropic,
\begin{equation}
  \operatorname{Cov}(\Delta x)=2A\Sigma_aA^\top+2\sigma^2I.
\end{equation}
Thus the leading eigenspace of paired differences estimates the acquisition subspace. A first-order biological projection is
\begin{equation}
  P_{\perp A}=I-AA^\top,\qquad x_{\mathrm{bio}}=P_{\perp A}x.
\end{equation}

\subsection{Geometric overlap penalty}
Let $B$ and $A$ be orthonormal bases with canonical overlap $\rho$. Removing the acquisition subspace retains approximately
\begin{equation}
  E_{\mathrm{bio,retained}}\approx1-\rho^2.
\end{equation}
Hence $\rho=0.90$ leaves only $0.19$ of the biological energy outside the acquisition subspace, while $\rho=0.95$ leaves $0.0975$. Any component in $\operatorname{span}(B)\cap\operatorname{span}(A)$ is not uniquely assignable without additional assumptions.

\begin{table}[t]
\centering
\caption{Geometric difficulty proxy.}
\small
\begin{tabular}{ccc}
\toprule
$\rho$ & $1-\rho^2$ & $(1-\rho^2)^{-2}$\\
\midrule
0.50 & 0.7500 & 1.78\\
0.75 & 0.4375 & 5.22\\
0.90 & 0.1900 & 27.70\\
0.95 & 0.0975 & 105.19\\
\bottomrule
\end{tabular}
\end{table}

\subsection{Finite-pair sample complexity}
Let $m=pn$ be the number of paired examples, $D$ the observed dimension, and $\gamma$ the eigengap separating acquisition directions from noise. A Davis--Kahan-style perturbation argument suggests
\begin{equation}
  \sin\Theta(\widehat A,A)
  \lesssim
  \frac{C}{\gamma}
  \sqrt{\frac{D+\log(1/\delta)}{pn}}.
\end{equation}
Requiring this estimation error to be smaller than the surviving biological margin gives the scaling law
\begin{equation}
  p_{\min}
  \gtrsim
  \frac{C^2(D+\log(1/\delta))}
  {n\gamma^2(1-\rho^2)^2}.
\end{equation}
This predicts diminishing returns in paired coverage, error decay proportional to $(pn)^{-1/2}$, and a sharp rise in required pairing as $\rho\rightarrow1$.

\subsection{Objective conflict in transport models}
For an operator trained with task, environment, reconstruction, transport, and consistency losses,
\begin{equation}
  \mathcal L=\mathcal L_{\mathrm{task}}+\lambda_e\mathcal L_{\mathrm{env}}+
  \lambda_r\mathcal L_{\mathrm{recon}}+\lambda_t\mathcal L_{\mathrm{transport}}+
  \lambda_c\mathcal L_{\mathrm{consistency}},
\end{equation}
the transport and consistency objectives may conflict. Their gradient cosine is
\begin{equation}
  G_{ct}=
  \frac{\nabla\mathcal L_c^\top\nabla\mathcal L_t}
  {\|\nabla\mathcal L_c\|\,\|\nabla\mathcal L_t\|}.
\end{equation}
A negative $G_{ct}$ indicates that reconstruction of acquisition-specific detail is opposing biological invariance.

\subsection{Imperfect interventions and wrong pairs}
If an intervention also changes biology, $z_b'=z_b+\eta_b$, then
\begin{equation}
  \operatorname{Cov}(\Delta x)=
  B\Sigma_{\eta_b}B^\top+2A\Sigma_aA^\top+2\sigma^2I.
\end{equation}
The ratio
\begin{equation}
  R=\frac{\lambda_{\max}(B\Sigma_{\eta_b}B^\top)}
  {\lambda_{\min}^{+}(2A\Sigma_aA^\top)}
\end{equation}
acts as an intervention-contamination measure: recovery should degrade as $R$ approaches one.

If a fraction $q$ of pairs are incorrect, random pair differences introduce biological variation. To first order,
\begin{equation}
  \Sigma_{\Delta}
  \approx2qB\Sigma_bB^\top+2A\Sigma_aA^\top+2\sigma^2I.
\end{equation}
Unlike finite-sample variance, this creates an asymptotic bias floor that more corrupted pairs cannot remove.

\subsection{Testable predictions}
The calculations predict: (i) minimum pairing grows roughly as $(1-\rho^2)^{-2}$; (ii) clean paired-subspace error decreases approximately as $(pn)^{-1/2}$; (iii) operator underperformance correlates with negative $G_{ct}$; and (iv) biology-changing interventions and wrong pairs create bias floors rather than merely slower convergence.
